% Conclusiones
\section{Conclusions}
\label{sec:conclusions}

We have adapted the block-based beam search BSG, originally designed for the single container
loading problem, to the palletization of air cargo into ULDs with non-cuboid shapes, a floor
edge and a substructure, weight capacities, fragile boxes, support constraints and CoG
constraints. The ULD constraints are handled in the generation and placement of the blocks, and
the CoG constraint in the value of the solutions and with a final
translation of the load, which is kept grouped by anchoring the free spaces at one corner of the
ULD. The ULDs are loaded sequentially, and the total volume of the ULDs is then reduced by
removing a ULD, inserting its boxes into the others and loading the rest into at most two smaller
ULDs.

On the 300 instances of \citet{Paquay2018b}, and with the same parameters in all the settings,
our algorithm improves on the RGS of \citet{Hessler2025} in both volume and weight distribution.
It uses 2.9--7.1\% less ULD volume in the original setting, 3.8--8.2\% less in the adapted one
and 3.4--7.2\% less with one ULD of each type, the smaller savings at a comparable running time,
and it also leaves less volume unloaded; at the same time, at most 1\% of its ULDs violate the
CoG constraint, against 7--11\% of the ULDs of the RGS. Most of the savings come from the reduction phase,
which looks at the whole solution; preferring smaller ULD types during the construction helps
much less.

These results show that the ideas that made block-based tree search successful in container
loading, namely building the load with blocks of boxes, evaluating them with the VCS function
and guiding a beam search with greedy completions, carry over to highly heterogeneous air
cargo, provided that the blocks and their placement respect the support and fragility
constraints and that the weight distribution is taken into account during the search.

Several directions remain open. A good lower bound would tell how far the remaining
solutions are from the optimum. The sequential framework could be replaced by a search that
decides the assignment of the boxes to the ULDs jointly, which matters most when the number of
ULDs is limited. Other practical constraints, such as load bearing capacities, the grouping of
the boxes by destination or the balance of the whole aircraft, could also be integrated into
the evaluation of the blocks and of the solutions.
